\section{高阶微商与高阶微分}\label{sec:4.4}

\subsection{高阶微商}

物体运动的瞬时速度 $v(t)$ 是路程 $s$ 关于时间 $t$ 的变化率，即
\[
v(t)=s'(t)=\frac{\mathrm{d}s}{\mathrm{d}t}.
\]
而物体运动的瞬时加速度 $a(t)$ 则是速度 $v$ 关于时间 $t$ 的变化率，即
\[
a(t)=v'(t)=(s'(t))',
\]
或
\[
a(t)=\frac{\mathrm{d}v}{\mathrm{d}t}=\frac{\mathrm{d}}{\mathrm{d}t}\left(\frac{\mathrm{d}s}{\mathrm{d}t}\right).
\]
由此产生了导函数的导数即高阶导数的概念．

一般地，设 $y=f(x)$ 在 $(a,b)$ 可导，则 $f'(x)$ 仍是 $(a,b)$ 上的函数．若 $f'(x)$ 也在 $(a,b)$ 可导，则称 $f'(x)$ 的微商 $(f'(x))'$ 为 $f(x)$ 的二阶微商（二阶导数），记为
\[
f''(x)\quad\text{或}\quad f^{(2)}(x)\quad\text{或}\quad\frac{\mathrm{d}^{2}y}{\mathrm{d}x^{2}}.
\]
类似地，可定义 $f''(x)$ 的微商为 $f(x)$ 的三阶微商（三阶导数），记为
\[
f'''(x)\quad\text{或}\quad f^{(3)}(x)\quad\text{或}\quad\frac{\mathrm{d}^{3}y}{\mathrm{d}x^{3}}.
\]
定义 $f^{(n-1)}(x)$ 的微商为 $f(x)$ 的 $n$ 阶微商（$n$ 阶导数），记为
\[
f^{(n)}(x)\quad\text{或}\quad\frac{\mathrm{d}^{n}y}{\mathrm{d}x^{n}}=\frac{\mathrm{d}}{\mathrm{d}x}\left(\frac{\mathrm{d}^{n-1}y}{\mathrm{d}x^{n-1}}\right).
\]
求函数的高阶微商，有时可用归纳法，得出一个一般的公式．

\begin{example}\label{ex:4-4-1}
设 $y=x^{n}$（$n$ 是正整数），求 $y$ 的各阶导数．

解
\[
\begin{aligned}
y'&=nx^{n-1},\\
y''&=n(n-1)x^{n-2},\\
y^{(n-1)}&=n(n-1)\cdots 2x,\\
y^{(n)}&=n(n-1)\cdots 2\cdot 1=n!,\\
y^{(k)}&=0\quad(k>n).
\end{aligned}
\]
\end{example}

\begin{example}\label{ex:4-4-2}
设 $y=\sin x$，求 $y^{(n)}$．

解 $y'=\cos x$，$y''=-\sin x$，$y'''=-\cos x$，$y^{(4)}=\sin x$，$\cdots$．

由此可见，每求四次导数将重复上述过程．我们知道 $\sin x$ 与 $\cos x$ 可以互相转换，因此，有可能获得一个统一的表达式．为此，作下面恒等变形
\[
\begin{aligned}
y'&=\cos x=\sin\left(x+\frac{\pi}{2}\right),\\
y''&=\cos\left(x+\frac{\pi}{2}\right)=\sin\left(x+\frac{\pi}{2}+\frac{\pi}{2}\right)=\sin\left(x+\frac{2\pi}{2}\right),\\
y'''&=\cos\left(x+\frac{2\pi}{2}\right)=\sin\left(x+\frac{2\pi}{2}+\frac{\pi}{2}\right)=\sin\left(x+\frac{3\pi}{2}\right).
\end{aligned}
\]
由此我们不难归纳出
\[
y^{(n)}=\sin\left(x+\frac{n\pi}{2}\right).
\]
对于 $y=\cos x$，可以用上面类似地方法求 $y^{(n)}$，也可以用上面的结论求 $y^{(n)}$，有
\[
y=\cos x=\sin\left(x+\frac{\pi}{2}\right),
\]
\[
y^{(n)}=\sin\left(x+\frac{\pi}{2}+\frac{n}{2}\pi\right)=\cos\left(x+\frac{n\pi}{2}\right).
\]
\end{example}

\begin{example}\label{ex:4-4-3}
设 $y=\arctan x$，求 $y^{(n)}$．

解
\[
y'=\frac{1}{1+x^{2}}=\frac{1}{1+\tan^{2}y}=\cos^{2}y;
\]
方程两边再对 $x$ 求导并注意 $y$ 是 $x$ 的函数，得
\[
\begin{aligned}
y''&=-2\cos y\sin y\cdot y'\\
&=-\sin 2y\cdot\cos^{2}y=\cos^{2}y\sin 2\left(y+\frac{\pi}{2}\right);
\end{aligned}
\]
\[
\begin{aligned}
y'''&=\left[-2\cos y\sin y\sin 2\left(y+\frac{\pi}{2}\right)+2\cos 2\left(y+\frac{\pi}{2}\right)\cos^{2}y\right]y'\\
&=2\cos^{3}y\left[\cos 2\left(y+\frac{\pi}{2}\right)\cos y-\sin y\sin 2\left(y+\frac{\pi}{2}\right)\right]\\
&=2\cos^{3}y\cos\left(3y+\frac{2}{2}\pi\right)\\
&=2\cos^{3}y\sin\left(3y+\frac{3}{2}\pi\right)\\
&=2\cos^{3}y\sin 3\left(y+\frac{1}{2}\pi\right);
\end{aligned}
\]
若
\[
y^{(n)}=(n-1)!\cos^{n}y\sin n\left(y+\frac{\pi}{2}\right),
\]
则
\[
\begin{aligned}
y^{(n+1)}&=(n-1)!\left[-n\cos^{n-1}y\sin y\sin n\left(y+\frac{\pi}{2}\right)+n\cos^{n}y\cos n\left(y+\frac{\pi}{2}\right)\right]y'\\
&=n!\cos^{n+1}y\left[\cos y\cos n\left(y+\frac{\pi}{2}\right)-\sin y\sin n\left(y+\frac{\pi}{2}\right)\right]\\
&=n!\cos^{n+1}y\cos\left[(n+1)y+\frac{n}{2}\pi\right]\\
&=n!\cos^{n+1}y\sin(n+1)\left(y+\frac{\pi}{2}\right).
\end{aligned}
\]
由数学归纳法得
\[
y^{(n)}=(n-1)!\cos^{n}y\sin n\left(y+\frac{\pi}{2}\right).
\]
\end{example}

从上两例可见，为了归纳出任意阶导数的表达式，恒等变形是一个重要的技巧．请读者设想，如果在例3中不用恒等变形
\[
\frac{1}{1+x^{2}}=\cos^{2}y,
\]
继续对 $\dfrac{1}{1+x^{2}}$ 求导会是怎样的情形．

高阶微商的运算法则也是重要的工具之一．显然，若 $u$、$v$ 都是 $x$ 的函数，则
\[
(u\pm v)^{(n)}=u^{(n)}\pm v^{(n)}.
\]
下面看乘积的高阶导数．设 $y=u\cdot v$，则
\[
y'=u'v+uv',
\]
\[
y''=u''v+2u'v'+uv''=\sum_{k=0}^{2}\mathrm{C}_{2}^{k}u^{(2-k)}v^{(k)},
\]
\[
y'''=u'''v+3u''v'+3u'v''+uv'''=\sum_{k=0}^{3}\mathrm{C}_{3}^{k}u^{(3-k)}v^{(k)},
\]
不难看出，
\[
y^{(n)}=\sum_{k=0}^{n}\mathrm{C}_{n}^{k}u^{(n-k)}v^{(k)},\quad n=1,2,\cdots.
\]
这里，函数的零阶导数理解为函数本身．下面用数学归纳法证明．事实上，设公式对 $n$ 成立，则
\[
\begin{aligned}
y^{(n+1)}&=\sum_{k=0}^{n}\mathrm{C}_{n}^{k}\left(u^{(n-k)}\cdot v^{(k)}\right)'\\
&=\sum_{k=0}^{n}\mathrm{C}_{n}^{k}\left[u^{(n-k+1)}v^{(k)}+u^{(n-k)}v^{(k+1)}\right]\\
&=\sum_{k=0}^{n}\mathrm{C}_{n}^{k}u^{(n+1-k)}v^{(k)}+\sum_{k=1}^{n+1}\mathrm{C}_{n}^{k-1}u^{(n+1-k)}v^{(k)}\\
&=\mathrm{C}_{n}^{0}u^{(n+1)}v^{(0)}+\sum_{k=1}^{n}\left(\mathrm{C}_{n}^{k}+\mathrm{C}_{n}^{k-1}\right)u^{(n+1-k)}v^{(k)}+\mathrm{C}_{n}^{n}u^{(0)}v^{(n+1)}\\
&=\sum_{k=0}^{n+1}\mathrm{C}_{n+1}^{k}u^{(n+1-k)}v^{(k)}.
\end{aligned}
\]
其中用到等式
\[
\mathrm{C}_{n}^{0}=\mathrm{C}_{n+1}^{0},\quad\mathrm{C}_{n}^{n}=\mathrm{C}_{n+1}^{n+1},\quad\mathrm{C}_{n}^{k}+\mathrm{C}_{n}^{k-1}=\mathrm{C}_{n+1}^{k}.
\]
由数学归纳法知公式对一切正整数 $n$ 成立．

乘积的 $n$ 阶微商公式
\[
y^{(n)}=\sum_{k=0}^{n}\mathrm{C}_{n}^{k}u^{(n-k)}v^{(k)}
\]
称为\emph{莱布尼茨公式}．显然，它与二项式展开公式非常相似，区别仅在于将乘幂次数换成了微商次数．

\begin{example}\label{ex:4-4-4}
设 $y=x^{3}\sin x$，求 $y^{(10)}$．

解
\[
\begin{aligned}
y^{(10)}&=x^{3}\sin\left(x+\frac{10}{2}\pi\right)+10\cdot 3x^{2}\sin\left(x+\frac{9}{2}\pi\right)+\frac{10\cdot 9}{2}\cdot 3\cdot 2x\sin\left(x+\frac{8}{2}\pi\right)+\frac{10\cdot 9\cdot 8}{3!}\cdot 3!\sin\left(x+\frac{7}{2}\pi\right)\\
&=x^{3}\sin\left(x+\frac{10}{2}\pi\right)+30x^{2}\sin\left(x+\frac{9}{2}\pi\right)+270x\sin\left(x+\frac{8}{2}\pi\right)+720\sin\left(x+\frac{7}{2}\pi\right).
\end{aligned}
\]
\end{example}

\begin{example}\label{ex:4-4-5}
设 $y=x^{m}\sin x$，其中 $m$ 是正整数，求 $y^{(n)}$．

解 由莱布尼茨公式，
\[
y^{(n)}=\sum_{k=0}^{n}\mathrm{C}_{n}^{k}\left(\sin x\right)^{(n-k)}\left(x^{m}\right)^{(k)}.
\]
由于 $k>m$ 时，$(x^{m})^{(k)}=0$，故当 $n>m$ 时，
\[
\begin{aligned}
y^{(n)}&=\sum_{k=0}^{m}\mathrm{C}_{n}^{k}\left(\sin x\right)^{(n-k)}\left(x^{m}\right)^{(k)}\\
&=\sum_{k=0}^{m}\mathrm{C}_{n}^{k}\sin\left(x+\frac{n-k}{2}\pi\right)m(m-1)\cdots(m-k+1)x^{m-k}\\
&=\sum_{k=0}^{m}\mathrm{C}_{n}^{k}\mathrm{C}_{m}^{k}k!\,x^{m-k}\sin\left(x+\frac{n-k}{2}\pi\right).
\end{aligned}
\]
\end{example}

\begin{example}\label{ex:4-4-6}
设 $y=\dfrac{1}{x(x-1)}$，求 $y^{(n)}$．

解
\[
y=\frac{1}{x(x-1)}=\frac{x-(x-1)}{x(x-1)}=\frac{1}{x-1}-\frac{1}{x}.
\]
易证 $\left(\dfrac{1}{x}\right)^{(n)}=(-1)^{n}\dfrac{n!}{x^{n+1}}$，故
\[
y^{(n)}=\left(\frac{1}{x-1}\right)^{(n)}-\left(\frac{1}{x}\right)^{(n)}=(-1)^{n}n!\left(\frac{1}{(x-1)^{n+1}}-\frac{1}{x^{n+1}}\right).
\]
注意，对 $y=\dfrac{1}{x(x-1)}=\dfrac{1}{x}\cdot\dfrac{1}{x-1}$ 用莱布尼茨公式当然可以，但显然是自找麻烦．
\end{example}

\subsection{高阶微分}

函数 $y=f(x)$ 的一阶微分是
\[
\mathrm{d}y=f'(x)\mathrm{d}x.
\]
其中 $x$ 和 $\mathrm{d}x$ 是两个独立的变量，现在把一阶微分视为 $x$ 的函数，如果它是可微的，则再求一次微分得
\[
\mathrm{d}(\mathrm{d}y)=\mathrm{d}(f'(x)\mathrm{d}x)=f''(x)(\mathrm{d}x)^{2}.
\]
上式称为函数 $y=f(x)$ 的二阶微分，记为 $\mathrm{d}^{2}y$．把 $(\mathrm{d}x)^{2}$ 记为 $\mathrm{d}x^{2}$，即有
\[
\mathrm{d}^{2}y=f''(x)\mathrm{d}x^{2}.
\]
注意 $\mathrm{d}x^{2}=(\mathrm{d}x)^{2}$ 是自变量微分的平方，不要把它与 $\mathrm{d}(x^{2})$ 及 $\mathrm{d}^{2}x$ 混淆．事实上，
\[
\mathrm{d}(x^{2})=2x\mathrm{d}x
\]
是函数 $y=x^{2}$ 的一阶微分，而 $\mathrm{d}^{2}x$ 则应理解为 $x$ 的二阶微分．

类似地，可以定义 $y=f(x)$ 的三阶微分
\[
\mathrm{d}^{3}y=f'''(x)\mathrm{d}x^{3}.
\]
一般地，$y=f(x)$ 的 $n$ 阶微分为
\[
\mathrm{d}^{n}y=f^{(n)}(x)\mathrm{d}x^{n},
\]
于是
\[
\frac{\mathrm{d}^{n}y}{\mathrm{d}x^{n}}=f^{(n)}(x).
\]
这正是 $n$ 阶导数符号 $\dfrac{\mathrm{d}^{n}y}{\mathrm{d}x^{n}}$ 的由来．

完全类似 $n=2$ 的情形，读者应搞清楚 $\mathrm{d}x^{n}$、$\mathrm{d}(x^{n})$ 和 $\mathrm{d}^{n}x$ 的差别．当 $n\geqslant 2$ 时，$n$ 阶微分统称为高阶微分．

同样可导出 $n$ 阶微分的运算法则．设 $u$、$v$ 都是 $x$ 的函数，则
\[
\mathrm{d}^{n}(u\pm v)=\mathrm{d}^{n}u\pm\mathrm{d}^{n}v,
\]
\[
\mathrm{d}^{n}(u\cdot v)=\sum_{k=0}^{n}\mathrm{C}_{n}^{k}\mathrm{d}^{n-k}u\,\mathrm{d}^{k}v,\quad n=1,2\cdots.
\]
一阶微分具有形式不变性，高阶微分是否也具有形式不变性呢？即当 $x$ 是中间变量时，公式 $\mathrm{d}^{n}y=f^{(n)}(x)\mathrm{d}x^{n}$ 是否仍成立呢？请看下例：

设 $y=e^{x}$，当 $x$ 是自变量时有
\[
\mathrm{d}^{2}y=e^{x}\mathrm{d}x^{2}.
\]
又若 $x=t^{2}$，则复合函数为 $y=e^{t^{2}}$，故
\[
\mathrm{d}^{2}y=\left(e^{t^{2}}\right)''\mathrm{d}t^{2}=\left(2e^{t^{2}}+4t^{2}e^{t^{2}}\right)\mathrm{d}t^{2},
\]
但
\[
e^{x}\mathrm{d}x^{2}=4t^{2}e^{t^{2}}\mathrm{d}t^{2}.
\]
可见当 $x$ 是中间变量时，$\mathrm{d}^{2}y=e^{x}\mathrm{d}x^{2}$ 不再成立，它少了一项 $2e^{t^{2}}\mathrm{d}t^{2}$．所以高阶微分不再具有形式不变性．

一般地若 $y=f(x)$，$x=g(t)$，由一阶微分形式不变性有
\[
\mathrm{d}y=f'(x)\mathrm{d}x.
\]
由于这时 $x$ 是中间变量，故 $\mathrm{d}x$ 和 $x$ 不再独立，它们都是自变量的函数，在求二阶微分时应该用乘积的微分法则，即
\[
\mathrm{d}^{2}y=f''(x)\mathrm{d}x^{2}+f'(x)\mathrm{d}^{2}x.
\]
与 $x$ 是自变量情形比较，它多了第二项，这就说明了高阶微分不具有形式不变性．因此，在带有高阶微分的等式中，不能随便使用变量代换．这是高阶微分与一阶微分的重要差别．

\begin{xiti}
\item 求下列函数在指定点的高阶导数：
\begin{ziti}
\item $f(x)=3x^{3}+4x^{2}-5x-9$，求 $f''(1)$，$f'''(1)$，$f^{(4)}(1)$；
\item $f(x)=\dfrac{x}{\sqrt{1+x^{2}}}$，求 $f''(0)$，$f''(1)$，$f''(-1)$；
\item $y=\dfrac{\arcsin x}{\sqrt{1-x^{2}}}$，求 $y^{(n)}(0)$．
\end{ziti}

\item 求下列函数的高阶导数：
\begin{ziti}
\item $y=x\ln x$，求 $y''$；
\item $y=e^{-x^{2}}$，求 $y''$；
\item $y=x^{2}e^{2x}$，求 $y^{(n)}$；
\item $y=x^{5}\cos x$，求 $y^{(50)}$；
\item $y=x^{3}\dfrac{e^{x}-e^{-x}}{2}$，求 $y^{(30)}$．
\end{ziti}

\item 求下列函数的 $n$ 阶导数：
\begin{ziti}
\item $y=a^{x}$；
\item $y=\ln x$．
\end{ziti}

\item 求下列函数的 $n$ 阶导数：
\begin{ziti}
\item $y=\dfrac{1}{x(1-2x)}$；
\item $y=\sin^{2}x$；
\item $y=\dfrac{1}{x^{2}-2x-8}$；
\item $y=\dfrac{e^{x}}{x}$；
\item $y=\ln\dfrac{x+2}{1-x}$；
\item $y=2^{x}\ln x$．
\end{ziti}

\item 设 $f(x)$ 的各阶导数存在，求 $y'$ 及 $y''$：
\begin{ziti}
\item $y=f(x^{2})$；
\item $y=f\left(\dfrac{1}{x}\right)$；
\item $y=f(e^{-x})$；
\item $y=f(\ln x)$；
\item $y=f(f(x))$．
\end{ziti}

\item 若 $f(x)=\begin{cases}e^{-\frac{1}{x^{2}}}, & x\neq0,\\ 0, & x=0,\end{cases}$ 证明 $f^{(n)}(0)=0$．

\item 求下列函数的二阶微分：
\begin{ziti}
\item $y=\dfrac{1}{\sqrt{x}}$；
\item $y=x\arctan x$；
\item $y=f(u)=e^{u}$，$u=\varphi(x)=x^{2}$．
\end{ziti}

\item 求下列函数的三阶微分：
\begin{ziti}
\item 设 $u(x)=\ln x$，$v(x)=e^{x}$，求 $\mathrm{d}^{3}(uv)$，$\mathrm{d}^{3}\left(\dfrac{u}{v}\right)$；
\item 设 $u(x)=e^{\frac{x}{2}}$，$v(x)=\cos 2x$，求 $\mathrm{d}^{3}(uv)$，$\mathrm{d}\left(\dfrac{u}{v}\right)$．
\end{ziti}

\item 求下列参数方程的二阶导数：
\begin{ziti}
\item $\begin{cases}x=2t-t^{2},\\ y=3t-t^{3};\end{cases}$
\item $\begin{cases}x=a\cos t,\\ y=a\sin t;\end{cases}$
\item $\begin{cases}x=a(t-\sin t),\\ y=a(1-\cos t);\end{cases}$
\item $\begin{cases}x=e^{t}\cos t,\\ y=e^{t}\sin t;\end{cases}$
\item $\begin{cases}x=a\cos^{3}t,\\ y=a\sin^{3}t;\end{cases}$
\item $\begin{cases}x=f'(t),\\ y=tf'(t)-f(t).\end{cases}$
\end{ziti}

\item 求下列隐函数的二阶导数 $\dfrac{\mathrm{d}^{2}y}{\mathrm{d}x^{2}}$：
\begin{ziti}
\item $e^{x+y}-xy=0$；
\item $x^{3}+y^{3}-3axy=0$；
\item $y^{2}+2\ln y-x^{4}=0$．
\end{ziti}

\item 设函数 $y=f(x)$ 在点 $x$ 二阶可导，且 $f'(x)\neq0$．若 $f(x)$ 存在反函数 $x=f^{-1}(y)$，试求 $(f^{-1})''(y)$．

\item 设 $y=c_{1}\sin x+c_{2}\cos x$，证明 $y$ 满足方程 $y''+y=0$．

\item 设 $y=\arctan x$．
\begin{ziti}
\item 证明 $y$ 满足方程 $(1+x^{2})y''+2xy'=0$；
\item 求 $y^{(n)}(0)$．
\end{ziti}

\item 设 $y=y(x)$ 存在反函数，且满足方程
\[
\frac{\mathrm{d}^{2}y}{\mathrm{d}x^{2}}+\left(\frac{\mathrm{d}y}{\mathrm{d}x}\right)^{3}=0.
\]
证明：反函数 $x=x(y)$ 满足 $\dfrac{\mathrm{d}^{2}x}{\mathrm{d}y^{2}}=1$，并由此求出一个 $y=y(x)$．
\end{xiti}
